\documentclass[11pt,a4paper]{article}
\usepackage[margin=23mm]{geometry}
\usepackage{amsmath,amssymb,amsthm}
\setlength{\emergencystretch}{2em}
\newtheorem{theorem}{Theorem}
\newcommand{\gen}[1]{\langle#1\rangle}
\newcommand{\Alt}{\operatorname{Alt}}
\title{An internal row gives a nonterminal separator at odd degree}
\author{Human Compression project: L044 continuation}
\date{}
\begin{document}\maketitle
Let $G=\gen{\rho_0,\ldots,\rho_{r-1}}\le\Alt(\Omega)$
be transitive, with a connected string C-marking and all
full maximal omissions intransitive. Assume $n=2r+1$ and
an actual transitive six-point head supported on
$D=\{1,\ldots,6\}$. Retain the literal hypotheses
\[
 \tau=\rho_3=(6\ b)\delta,\quad
 \operatorname{supp}(\delta)\subseteq\Omega\setminus(D\cup\{b\}),
 \qquad R=G_{\ge3}=\Alt(\Omega\setminus\{1,\ldots,5\}).
\]
Put $\sigma=\rho_4$, $H=G_{\ge5}$ and $K=G_{>5}$.
Assume $\sigma$ fixes $D$, while $H$ fixes $D\cup\{b\}$.
As in the acquired pure-overlap proof, $x=b^\sigma\ne b$,
$\delta x=x$, and $\gen{\tau,\sigma}$ contains a pure
three-cycle on $\{6,b,x\}$. These follow from the displayed
actual natural $R$ and full IP, without a rank induction.

\begin{theorem}
Let $O$ be an actual $H$-orbit of size $2q$, $q\ge3$,
which $\delta$ preserves and on which it acts nontrivially. Suppose $K$ is
intransitive on $O$. Let $F\subseteq\{6,\ldots,r-1\}$
have size $q-1$, with an ambient-eligible edge of every
colour in $F$ contained in $O$. Then $F$ is a consecutive
interval, and its last label $z$ satisfies
\[
 6\le z<r-1,\qquad
                  |\operatorname{Fix}(G_{>z})|\ge2z+3.
\]
In particular it supplies the strict nonterminal separator
needed at this internal-orbit equality.
\end{theorem}
\begin{proof}
The involution $\delta$ centralizes $H$, as $H$ fixes
$6,b$ and centralizes $\tau$. It is free on the transitive
orbit $O$. Moreover $\tau\in G_j$ for each $j\in F$,
and its restriction to $O$ is $\delta|_O$.
The complete acquired internal-pair theorem therefore applies
with $k=5$. Write $a=\min F$ and $z=\max F$.
It gives
\[
 F=\{a,\ldots,a+q-2\},\qquad z=a+q-2,
\]
and $K|_O$ is the diagonal natural $S_q$ on two rows.
Every tail mark outside $F\cup\{5\}$ fixes $O$ pointwise.
The acquired prefix-row growth and its additional-fixed-point
lemma, using the seven actual $H$-fixed points $D\cup\{b\}$,
give the stronger estimate
\[
              |\operatorname{Fix}(G_{>z})|\ge(a-3)q+8.       \tag{1}
\]
Both acquired lemmas hold before the last label as well as
at it; they do not assume the critical even-degree equation.
Writing $a=6+e$, $e\ge0$, the difference between (1)
and $2z+3$ is
\[
                          q-3+e(q-2)\ge0.
\]
This proves the required fixed-point bound.

Suppose that $z=r-1$. Then every point is fixed by
$G_{>z}$ and $n=2r+1=2z+3$. Equality in the last inequality
is necessary, so $q=3$, $e=0$, $a=6$, $z=7$, $r=8$
and $n=17$. The two $K=\gen{\rho_6,\rho_7}$ rows
of $O$ have three points each. The complement of
$D\cup\{b\}\cup O$ is a four-point set $V$, preserved
by $H$ and $K$. The point $x$ lies in $V$: it is outside
$D\cup\{b\}$ and is fixed by $\delta$, which is free on
$O$. Since $K$ fixes $b$ and commutes with $\sigma$, it
fixes $x$ as well.

The group $R=\gen{\tau,\sigma,H}$ is transitive on its
natural support. Since $\tau$ and $H$ preserve $O$,
$\sigma$ cannot preserve $O$. It centralizes $K$, so it
must send one of the two three-point $K$-orbits in $O$
to a three-point $K$-orbit outside $O$. All points of
$D\cup\{b,x\}$ are fixed by $K$. Therefore that new
orbit is exactly $W=V\setminus\{x\}$.
The two involutions $\rho_6,\rho_7$ act transitively on
$W$, and hence as two distinct transpositions generating
natural $S_3$. They likewise act as the two path
transpositions on each row of $O$. Their product has order
three on all three rows and they fix the remaining points.
Thus $K$ itself is diagonal $S_3$ on these three natural rows.

Its centralizer on their union is exactly $S_3$ permuting
the rows: each natural row has trivial centralizer and the
equivariant bijection between any two rows is unique.
The involution $\delta$ swaps the two rows of $O$, since
it is free on $O$; it fixes $W$ pointwise, since it preserves
the three-point $K$-orbit $W$ and centralizes its transitive
action. It also fixes $D,b,x$.
The involution $\sigma$ swaps $b,x$, fixes $D$, and must
exchange $W$ with one of the two rows of $O$ because it
does not preserve $O$. Hence their exterior row actions
are two distinct transpositions. On the nine-point row
union the product $\tau\sigma=\delta\sigma$ has order
three. On $\{6,b,x\}$ it is also a three-cycle, and
it fixes the other five head points.

Thus $\gen{\tau,\sigma}$ is a dihedral group whose
nonidentity rotations have nontrivial restrictions on both
sets. It cannot contain the acquired pure three-cycle on
$\{6,b,x\}$. This contradiction excludes $z=r-1$ and
finishes the proof.
\end{proof}

\paragraph{Attribution and boundary.}
The internal-pair equality, row-growth estimate and extra fixed
point are complete acquired proofs. The new argument eliminates
their only possible terminal equality at $n=2r+1$, using three
actual natural $S_3$ rows and the pure-cycle intersection.
No quotient IP or new computation is used. The existence of
such a saturated nonfirst internal orbit remains a hypothesis;
other orbit deficits and pairs split across different orbits
are still needed for the full separator refinement. No
whole-original, formal or question credit is claimed.
\end{document}
